%-----------------------------------------------------------------------------%
\chapter{\babSatu}
\label{Pendahuluan}
%-----------------------------------------------------------------------------%
Autism Spectrum Disorder (ASD) merupakan gangguan neurologis yang memengaruhi komunikasi, interaksi sosial, dan perilaku individu, sehingga proses identifikasi dini memiliki peranan penting dalam penanganannya. Analisis sinyal electroencephalography (EEG) menggunakan metode machine learning mulai banyak diteliti sebagai pendekatan dalam proses klasifikasi ASD secara lebih objektif. Meskipun demikian, pengembangan model klasifikasi berbasis EEG masih menghadapi tantangan berupa keterbatasan dataset EEG yang tersedia dan tingginya variabilitas sinyal antar individu maupun antar sesi perekaman. Karakteristik tersebut menyebabkan terhambatnya proses pengembangan model klasifikasi EEG yang mampu memberikan performa secara konsisten pada data dengan karakteristik yang berbeda. Salah satu pendekatan yang mulai banyak digunakan untuk mengatasi permasalahan tersebut adalah teknik \textit{data augmentation}, yaitu proses penambahan variasi data \textit{training} secara artifisial, sehingga diharapkan dapat meningkatkan performa dan kemampuan generalisasi model.

Bab ini akan dimulai dengan penjelasan Latar Belakang yang memberikan konteks penting terkait penelitian ini, diikuti oleh Rumusan Masalah dan Tujuan Penelitian yang merinci fokus utama studi. Selanjutnya, akan dibahas Manfaat Penelitian dan Ruang Lingkup Penelitian untuk menjelaskan kontribusi serta batasan penelitian ini. Terakhir, Sistematika Penulisan akan menguraikan struktur keseluruhan laporan penelitian.

%-----------------------------------------------------------------------------%
\section{Latar Belakang}
\label{sec:latar_belakang}
%-----------------------------------------------------------------------------%

Autism Spectrum Disorder (ASD) merupakan gangguan neurologis yang ditandai oleh hambatan dalam komunikasi dan interaksi sosial, serta pola perilaku yang terbatas dan berulang \cite{nih_autism}, \cite{dsm5}. Identifikasi dan intervensi dini menjadi faktor penting dalam membantu perkembangan individu dengan ASD, namun proses diagnosis masih menghadapi berbagai tantangan karena karakteristik gejala yang heterogen dan sering kali menyerupai gangguan perkembangan lainnya \cite{autism_diagnosis_challenges} \cite{10423494}. Kondisi tersebut mendorong berkembangnya pendekatan diagnosis berbasis biomarker yang lebih objektif dan konsisten.

Salah satu pendekatan yang banyak diteliti adalah analisis electroencephalography (EEG) menggunakan metode machine learning. Dibandingkan dengan metode lainnya, EEG memiliki keunggulan karena bersifat non-invasif, relatif lebih murah, dan memiliki resolusi temporal yang tinggi \cite{brain_imaging_asd_review} \cite{eeg_dataset_asd}. Berbagai penelitian menunjukkan bahwa model machine learning mampu mengidentifikasi pola aktivitas otak yang berkaitan dengan ASD melalui sinyal EEG dengan performa klasifikasi yang cukup baik \cite{eeg_autism_review}. Penelitian rujukan \cite{paper_ref} mengusulkan pendekatan klasifikasi ASD berbasis EEG menggunakan kombinasi seleksi fitur minimum Redundancy Maximum Relevance (mRMR) dan algoritma XGBoost yang menunjukkan performa klasifikasi yang baik dan konsisten.

Meskipun demikian, pengembangan model klasifikasi EEG masih menghadapi berbagai tantangan, terutama terkait keterbatasan dataset EEG yang tersedia secara publik dan telah memiliki label klinis \cite{dl_eeg_review}. Proses akuisisi dan pelabelan data EEG umumnya membutuhkan biaya, waktu, serta keterlibatan tenaga ahli, sehingga jumlah data yang tersedia sering kali terbatas. Keterbatasan dataset ini dapat meningkatkan risiko overfitting pada proses training model machine learning dan membatasi kemampuan model dalam melakukan generalisasi terhadap data baru \cite{eeg_limited_data}. Selain itu, sinyal EEG juga rentan terhadap berbagai artefak, seperti gerakan otot, kedipan mata, maupun gangguan pada perangkat perekaman yang dapat memengaruhi kualitas sinyal \cite{artifact_removal_eeg}. Keberadaan artefak tersebut menyebabkan proses analisis EEG menjadi lebih sulit karena sinyal yang diperoleh tidak selalu merepresentasikan aktivitas otak yang sebenarnya.

Salah satu pendekatan yang mulai banyak diteliti untuk mengatasi keterbatasan dataset EEG adalah penggunaan teknik data augmentation. Teknik ini bertujuan untuk meningkatkan jumlah dan variasi data training secara artifisial sehingga model dapat mempelajari representasi fitur yang lebih beragam. Beberapa penelitian menunjukkan bahwa penerapan data augmentation pada data EEG dapat membantu meningkatkan performa model klasifikasi dan mengurangi risiko overfitting pada kondisi dataset yang terbatas \cite{eeg_augmentation_review}.

Penerapan augmentasi pada data EEG memiliki tantangan tersendiri karena sinyal EEG merupakan data \textit{time series} multikanal yang bersifat non-stasioner dan memiliki rasio \textit{signal-to-noise} yang rendah \cite{eeg_feature_review}. Berbeda dengan berbagai jenis data multimedia lain seperti citra, audio, atau video yang umumnya memiliki pola representasi yang lebih stabil, manipulasi pada data EEG perlu dilakukan secara hati-hati karena perubahan pada sinyal dapat memengaruhi representasi aktivitas otak yang dianalisis. Pendekatan augmentasi data konvensional pada data \textit{time series} umumnya dilakukan melalui transformasi sederhana seperti penambahan \textit{noise}, pergeseran waktu, atau manipulasi amplitudo. Namun, pada data EEG, transformasi tersebut berpotensi mengubah karakteristik temporal dan spasial yang digunakan dalam proses analisis maupun klasifikasi sinyal \cite{gan_eeg_emotion}.

Seiring dengan perkembangan deep learning, pendekatan berbasis model generatif seperti Variational Autoencoder (VAE) mulai digunakan untuk menghasilkan data EEG sintetis yang lebih representatif. Model VAE mempelajari distribusi laten dari data EEG sehingga mampu menghasilkan sampel baru yang berbeda namun tetap mempertahankan karakteristik statistik dari data asli. Pendekatan generatif memiliki potensi untuk menghasilkan variasi data yang lebih realistis dibandingkan metode augmentasi konvensional \cite{gan_eeg_emotion}. Meskipun demikian, pemanfaatan data sintetis berbasis model generatif pada klasifikasi EEG, khususnya dalam konteks Autism Spectrum Disorder (ASD), masih memerlukan kajian lebih lanjut untuk memahami pengaruhnya terhadap performa dan kemampuan generalisasi model klasifikasi.

Berdasarkan permasalahan tersebut, penelitian ini bertujuan untuk mengeksplorasi penggunaan teknik data augmentation berbasis Variational Autoencoder (VAE) untuk menghasilkan data EEG sintetis dan menganalisis pengaruhnya terhadap performa model klasifikasi ASD berbasis EEG. 
\section{Rumusan Masalah}
\label{sec:rumusan_masalah}
%-----------------------------------------------------------------------------%
Berdasarkan latar belakang yang telah dijelaskan, permasalahan utama dalam penelitian
ini dirumuskan sebagai berikut:
\begin{enumerate}
    \item Bagaimana pendekatan yang dapat digunakan untuk mengatasi keterbatasan jumlah dan variasi data pada sinyal EEG dalam konteks klasifikasi Autism Spectrum Disorder (ASD)?
    
    \item Sejauh mana data sintetis yang dihasilkan melalui proses augmentasi mampu mempertahankan karakteristik dan distribusi sinyal EEG asli?
    
    \item Bagaimana pengaruh penggunaan data hasil augmentasi terhadap performa dan kemampuan generalisasi model klasifikasi EEG dalam mendeteksi pola yang berkaitan dengan Autism Spectrum Disorder (ASD)?
\end{enumerate}

%-----------------------------------------------------------------------------%
\section{Tujuan Penelitian}
\label{sec:tujuan_penelitian}
%-----------------------------------------------------------------------------%

Selaras dengan rumusan masalah yang telah diuraikan, tujuan dari penelitian ini mencakup sebagai berikut:

\begin{enumerate}
    \item Mengkaji dan mengidentifikasi pendekatan yang dapat digunakan untuk mengatasi keterbatasan jumlah dan variasi data pada sinyal EEG dalam konteks klasifikasi Autism Spectrum Disorder (ASD).

    \item Mengevaluasi sejauh mana data sintetis yang dihasilkan melalui proses augmentasi mampu mempertahankan karakteristik dan distribusi sinyal EEG asli.

    \item Menganalisis pengaruh penggunaan data hasil augmentasi terhadap performa serta kemampuan generalisasi model klasifikasi EEG dalam mendeteksi pola yang berkaitan dengan Autism Spectrum Disorder (ASD).
\end{enumerate}

%-----------------------------------------------------------------------------%
\section{Manfaat Penelitian}
\label{sec:manfaat_penelitian}
%-----------------------------------------------------------------------------%
Berdasarkan latar belakang, rumusan masalah, dan tujuan penelitian yang telah diuraikan, penelitian ini diharapkan memberikan manfaat sebagai berikut:

\begin{enumerate}
    \item \textbf{Bagi Peneliti dan Pengembangan Metode:} \\
    Penelitian ini diharapkan dapat menjadi referensi dalam pengembangan dan pemanfaatan teknik \textit{data augmentation} pada sinyal EEG, serta membuka peluang bagi pengembangan metode klasifikasi dan diagnosis di masa depan yang mampu memanfaatkan data dalam jumlah yang lebih besar dan beragam, sehingga meningkatkan robustness dan kemampuan generalisasi model.

    \item \textbf{Bagi Tenaga Medis dan Profesional Kesehatan:} \\
    Penelitian ini diharapkan dapat memberikan gambaran mengenai potensi pemanfaatan \textit{machine learning} dalam analisis data EEG pada studi Autism Spectrum Disorder (ASD), maupun untuk  sehingga dapat mendukung pemahaman dan penelitian lanjutan terkait pendekatan berbasis data dalam konteks klinis.
\end{enumerate}

%-----------------------------------------------------------------------------%
\section{Ruang Lingkup Penelitian}
\label{sec:ruang_lingkup}
%-----------------------------------------------------------------------------%
Agar penelitian ini tetap terfokus, maka ruang lingkup penelitian ditetapkan sebagai berikut:
\begin{enumerate}

\item \textbf{Dataset Penelitian:} Penelitian ini menggunakan dataset EEG yang sama dengan dataset pada penelitian rujukan \cite{paper_ref}\cite{eeg_dataset_asd}. Dataset tersebut digunakan sebagai dasar dalam proses pelatihan model augmentasi maupun model klasifikasi sehingga hasil penelitian dapat dibandingkan secara relevan dengan penelitian sebelumnya.

\item \textbf{Metode Preprocessing Data:} Tahapan \textit{preprocessing} data EEG mengikuti prosedur yang digunakan pada penelitian rujukan \cite{paper_ref}. Proses ini mencakup tahapan pengolahan data yang diperlukan sebelum proses seleksi fitur, augmentasi data, dan klasifikasi dilakukan.

\item \textbf{Metode Augmentasi dan Klasifikasi:} Penelitian ini menerapkan metode \textit{data augmentation} berbasis \textit{Variational Autoencoder} (VAE) untuk menghasilkan data EEG sintetis. Selanjutnya, proses klasifikasi dilakukan menggunakan pendekatan yang mengacu pada penelitian rujukan \cite{paper_ref}, yaitu kombinasi seleksi fitur \textit{mRMR} dan model klasifikasi \textit{XGBoost}.

\item \textbf{Metode Evaluasi:} Evaluasi penelitian dilakukan dengan membandingkan performa model klasifikasi pada data tanpa augmentasi dan dengan augmentasi. Penilaian kualitas data sintetis serta performa klasifikasi mengacu pada metode evaluasi yang digunakan dalam penelitian rujukan \cite{eeg_gan}.

\item \textbf{Hasil Akhir Penelitian:} Penelitian ini menghasilkan model \textit{data augmentation} EEG berbasis \textit{Variational Autoencoder} (VAE) yang dapat digunakan untuk menghasilkan data EEG sintetis guna mendukung peningkatan performa model klasifikasi.


\item \textbf{Batasan Penelitian:} Penelitian difokuskan pada analisis pengaruh penerapan \textit{data augmentation} EEG terhadap kualitas performa model klasifikasi. Penelitian ini tidak membahas aspek perangkat keras EEG, proses akuisisi data asli, maupun tahapan penghilangan \textit{noise} dan \textit{artifact removal} di luar prosedur \textit{preprocessing} yang telah ditetapkan pada penelitian rujukan.
\end{enumerate}
%-----------------------------------------------------------------------------%
\section{Sistematika Penulisan}
\label{sec:sistematika}
%-----------------------------------------------------------------------------%
Laporan tesis ini disusun secara sistematis ke dalam lima bab utama, dengan rincian se-
bagai berikut:

\begin{enumerate}
    \item \textbf{Bab 1: Pendahuluan.} 
    Bab ini menyajikan gambaran umum penelitian, meliputi latar belakang yang menjelaskan pentingnya analisis sinyal EEG dalam mendukung identifikasi Autism Spectrum Disorder (ASD), rumusan masalah yang menjadi dasar penelitian, tujuan yang ingin dicapai, manfaat penelitian, ruang lingkup penelitian, serta sistematika penulisan laporan ini.
    
    \item \textbf{Bab 2: Tinjauan Pustaka.} 
    Bab ini menguraikan landasan teoretis yang mendukung penelitian, meliputi konsep dasar sinyal EEG, metode ekstraksi dan seleksi fitur pada data EEG, metode \textit{Minimum Redundancy Maximum Relevance} (mRMR), implementasi \textit{machine learning}  dalam klasifikasi berdasarkan EEG, metode augmentasi data EEG, model \textit{Variational AutoEncoder}, serta kajian terhadap penelitian-penelitian sebelumnya yang bersifat relevan.
    
    \item \textbf{Bab 3: Metodologi Penelitian.} 
    Bab ini menjelaskan metode yang digunakan dalam penelitian, mencakup deskripsi dataset EEG yang digunakan, tahapan pra-pemrosesan data, penerapan teknik \textit{data augmentation}, proses seleksi fitur menggunakan metode \textit{mRMR}, arsitektur model augmentasi berbasis \textit{VAE}, serta konfigurasi \textit{training} model klasifikasi menggunakan \textit{XGBoost}.
    
    \item \textbf{Bab 4: Hasil dan Analisis.} 
    Bab ini memaparkan hasil eksperimen yang diperoleh dari proses \textit{training} dan pengujian model. Analisis dilakukan untuk membandingkan performa model pada skenario dengan dan tanpa penerapan \textit{data augmentasi}, serta mengevaluasi pengaruhnya terhadap kemampuan model dalam melakukan klasifikasi data EEG.
    
    \item \textbf{Bab 5: Penutup.} 
    Bab terakhir ini berisi kesimpulan dari hasil penelitian yang telah dilakukan serta saran untuk pengembangan penelitian selanjutnya terkait analisis dan klasifikasi data EEG menggunakan metode \textit{machine learning} .
\end{enumerate}